\section{Wnioski}
\label{sec:wnioski}

\subsection{Synteza i kluczowe odkrycie}

Najbardziej zaskakującym wynikiem projektu jest \textbf{nieskuteczność
RAVE w Breakthrough 8$\times$8}. Pokazuje to, że standardowe techniki rozwoju
MCTS są domain-specific: założenie ,,dobrego ruchu zagranego gdziekolwiek''
(kluczowe dla AMAF) działa świetnie w Go, lecz załamuje się w grze, gdzie
te same fizyczne ruchy mają radykalnie różną wartość w zależności od pozycji.
Co istotne, wynik ten nie jest artefaktem implementacji -- poprawność szkieletu
RAVE potwierdziliśmy testem $\beta \to 0$ (odtworzenie siły UCT), więc spadek
pochodzi z samej heurystyki AMAF.

Drugim ważnym wnioskiem jest \textbf{przygniatająca przewaga prostej heurystyki
alpha-beta} dla tej klasy gier. Heurystyka oparta na bogatej funkcji ewaluacji
(wykładnicze zaawansowanie, groźba przebicia, obrońcy, centralność, swobodna
ścieżka, izolacja, mobilność i materiał -- Algorytm~1) z głębokością 5 i
porządkowaniem ruchów dominuje MCTS w całym praktycznym zakresie iteracji
(94\% nawet przy 100\,000 iteracji UCT), grając przy tym rzędy wielkości szybciej.

Trzecim, praktycznym odkryciem jest \textbf{korzyść z niskiej stałej eksploracji}
(H5): w taktycznym Breakthrough $c = 0{,}5$ wyraźnie przewyższa kanoniczne
$c = \sqrt{2}$, co sugeruje, że gra premiuje eksploatację (głębsze drążenie
wymuszonych wariantów) nad eksploracją.

\subsection{Propozycje dalszych prac}

\begin{itemize}[nosep]
    \item \textbf{Polityka rolloutowa oparta na heurystyce.} Jednym z najprostszych
          rozszerzeń byłaby zamiana czysto losowych symulacji na epsilon-greedy
          po heurystyce (np. ulubione ruchy z funkcji ewaluacji). Powinno to
          drastycznie poprawić wszystkich agentów MCTS i może zmienić
          rezultat H1 (RAVE z heurystycznymi rolloutami).
    \item \textbf{MAST i NAST.} Alternatywne techniki uczące się polityki
          symulacyjnej -- Move-Average Sampling Technique \cite{finnsson2008simulation}
          -- mogłyby przewyższyć heurystykę w niskim/średnim zakresie
          iteracji. To naturalna kontynuacja H2.
    \item \textbf{Krzyżowanie z deep RL.} Implementacja prostej sieci ewaluacyjnej
          (jak w AlphaZero, lecz z mniejszą architekturą) i jej porównanie
          z heurystyką alpha-beta.
\end{itemize}
